\section{Limitations}
\label{sec:limitations}

These limits follow from the design and from what went wrong on the way. None
of them depends on a held-out outcome.

\paragraph{One simulator.}
% source: upstream-99p/docs/CITATION_AUDIT.md:235 (Fail2Drive runs on CARLA 0.9.15)
% source: docs/prd-discovery-sprint.md:30-31 (recovery of known groups as a proxy)
Every route was driven in CARLA, on Fail2Drive's routes. The gaps are scenario
types that the benchmark's authors built into that simulator. A method that
recovers them shows that it can find failure groups of this kind in this
simulator. It does not show that it finds the groups a real fleet meets, and
recovery of known groups stands in for discovering unknown ones.

\paragraph{Two to three driving models.}
% source: docs/prd-discovery-sprint.md:39-41 (model roles); docs/paper/data.tex:63-67 (pool table)
% source: docs/prereg-fail2drive.md:238-241 (round 2 makes TransFuser++ a second dev pool)
% source: docs/prd-discovery-sprint.md:45-47 (TransFuser++ tests across models, not types; third model unseen)
Methods were developed on SimLingo. TransFuser++ was held out for round 1 and
became a second dev pool in round 2. UniAD is the only pool that no method has
been tuned on. TransFuser++'s gap types are a subset of SimLingo's, so it tests
whether a method carries across models, not across failure types. A claim that
a method transfers rests on very few models, and on one model that no method
has seen.

\paragraph{Small gap categories.}
% measured: pairs per type from results/h21/h21b/manifest.json (pair with -sNN stripped):
%   14 types have 5 pairs and 3 types have 10; 15 or 30 slice routes per type over seeds 24-26
% design: 100 pairs, 17 types (docs/prereg-fail2drive.md:13; upstream-99p/docs/CITATION_AUDIT.md:235)
% design: gap floor of 30 slice runs, six test seeds, every pair in each (docs/prereg-fail2drive.md:24-25, :31)
% measured x design: at most 30 or 60 failing routes = 5 or 10 pairs x 6 seeds, one slice route each
Fail2Drive has 100 pairs over 17 types. In the H21b manifest, 14 types have 5
pairs and the other 3 have 10, so a type has 15 or 30 slice routes over the
three seeds. On a six-seed test pool, a 5-pair type has 30 slice runs, which is
exactly the gap rule's floor if every route completes. A gap's failing routes
can be no more than its slice routes, so on a six-seed pool a gap holds at most
30 or 60 failing routes.
% source: docs/WHAT_BROKE.md:160-162 (M1's group-size floor against the gap sizes)
% source: docs/prd-discovery-sprint.md:83 (primary score is the mean over gaps of the best Jaccard)
If every axis a method emits holds more routes than a gap has failing routes,
its Jaccard with that gap is capped below 1. This happened once already, to the frozen M1 (\S\ref{sec:whatbroke}). The
primary score is a mean over a pool's few gaps, so it moves a long way when one
gap's best axis changes.

\paragraph{Many tries on the same dev pools.}
% source: docs/prd-discovery-sprint.md:89-93; docs/prereg-fail2drive.md:153 (dev p values are exploratory; the held-out score is the result)
% measured: 232 methods tried (docs/research/experiments/f2d-dev-simlingo/VARIANTS.md:1495)
A dev $p$ value read after many tries marks a candidate, not a result. The try
count stands at 232 methods, and every dev table states the running count
beside its $p$ values. The design treats only the held-out score of a frozen
finalist as a result.
% source: docs/research/experiments/f2d-dev-simlingo/VARIANTS.md:125, :360 (H21 can rank and kill, not make a candidate)
% source: CONFLICTS.md:847 (C42: round 2's search stops with no new finalist)
% source: VARIANTS.md:1144, :1196, :1304-1305, :1356, :1391, :1430 (each j16-j27 batch: cannot enter the UniAD read)
% source: docs/paper/data.tex:96-97 (H21b: same routes, seeds 25 and 26)
H21's categories had already been read by H22 to H25, so H21 can rank and
kill methods but cannot make a candidate on its own. H21b drives the same
routes on two more seeds. The j rows run
after the search stopped (j16--j27) are exploratory and cannot enter the UniAD
read.

\paragraph{Window position.}
% source: docs/research/experiments/f2d-dev-simlingo/VARIANTS.md:108 (m19); docs/paper/whatbroke.tex:38-45
% source: docs/research/experiments/f2d-dev-simlingo/VARIANTS.md:1443-1446 (every cached window ends at the failure or the route end)
% source: docs/research/experiments/f2d-dev-simlingo/VARIANTS.md:1458-1459 (failing-route clusters and failure step; diagnostic, numbers not used)
% source: docs/paper/methods.tex:106-118 (m23 and m40 settings, no position removal)
An approach window ends at the failure on a failing route and at the route's
end on a passing one. Window features therefore carry where in the route the
window sits, and m19's probes read that position instead of anything about the
failure (\S\ref{sec:whatbroke}). Among failing routes, too, windows sit at
different points, because failures come at different times. A cluster of
failing routes can therefore follow when they failed rather than why. The finalists m23 and m40 do not
remove position.
% source: docs/research/experiments/f2d-dev-simlingo/j26_attr.py:9-15, :98-108 (iterative linear projection)
% design: R^2 tolerance 0.05 (j26_attr.py TOLERANCE)
% measured: boosted-tree regressor R^2 0.588 on window-end step after removal (VARIANTS.md:1454-1457, results/h21/j26/position-check.json)
j26 removes position only linearly. It projects out ridge-regression
directions until no ridge fit predicts position with a cross-fitted $R^2$ above
0.05. Among failing routes, a boosted-tree regressor still predicted window
position from the features that remained, so a nonlinear probe such as j26's
may still read it. How much j26's probe does was not measured.

\paragraph{The H21b pair bug.}
% source: docs/WHAT_BROKE.md:19-28; VARIANTS.md:1065-1067 (intervals and p values too narrow)
% source: docs/research/experiments/f2d-dev-simlingo/sprint_score.py:106-108 (pair taken from the manifest as written)
The pair bootstrap is only as sound as the pair field it reads. When the three
H21 seeds were pooled into H21b, the pool builder copied each route's
\code{base\_id} into the manifest's \code{pair}. That id carries the seed, so
the scorer resampled pair-and-seed units instead of pairs, against A1. The
intervals it recorded were too narrow and the $p$ values too small. The point
scores were not affected. An independent scorer written from the pre-registration text
caught it. Every later H21b score uses a manifest copy with the seed suffix
stripped, and the UniAD test manifest was checked and has no suffix. Nothing in
the scorer refuses a pair field that carries a seed, so the check rests on the
pool builder (\S\ref{sec:h21b}).

\paragraph{Label hiding is mostly a convention.}
% source: docs/prereg-fail2drive.md:68-71 (the plan: shards exported without label fields, salted-hash renaming)
% source: adapters/simlingo/export_activations.py:146-161; src/harness/shards.py:7-14 (shard sidecars keep route_id, base_id, scenario, practice_gaps, role)
% source: docs/research/experiments/f2d-dev-simlingo/cache_pool.py:13-14, :32-33, :51, :68, :75; frames_pool.py:8-10, :134, :186 (salted row order, no scenario field, raw route_id kept)
% source: scripts/check_sealed_imports.py:1-36 (guarded directories and forbidden patterns)
% source: CONFLICTS.md:693-699 (C33: CI off, so the guard runs only by hand)
Two things enforce label hiding in code. The caches the dev methods read write
no scenario, town or route-config field and list routes in a salted-hash
order, but they keep each route's raw id, and the exported shards still carry
the type and role. The script
\code{scripts/check\_sealed\_imports.py} fails if a Python file under
\code{discovery/}, \code{baselines/} or \code{gate/}, at the top level or
under \code{src/}, imports the sealed key or names the sealed folder, the key
file or the salt. The dev method scripts live in
\code{docs/research/experiments/}, outside the guard's scope. Continuous
integration is off, so the guard runs only when someone runs it before a
commit.
% measured: results/h21/h21b/manifest.json rows carry `scenario` and `practice_gaps` (the Fail2Drive type); file not tracked in git (read in this session)
% source: docs/research/experiments/f2d-dev-simlingo/VARIANTS.md:1467-1468 (j26 row: never reads routes.jsonl, the manifest or route types)
% source: docs/WHAT_BROKE.md:152 (H25: id order sorted by category; answer mapping in the frames' parent folder)
The dev manifests give each route's Fail2Drive type, and they sit on the same
disk as the method scripts and the pools those scripts read. That a method never
opens them is a rule written into rows of \code{VARIANTS.md} such as j26's, and
the guard does not check it. The H25 review
showed how such hiding can fail. Route ids that sorted by category, and an
answer mapping one folder above the frames, were caught by a review before the
lock, not by code.
